\honest{Biases: An Introduction}{Rob Bensinger}{2015}

Reach into an urn of seventy white balls and thirty red ones, and take out ten. By chance you may draw three red and guess the mix correctly, or four and guess wrong. Such random errors are right on average and shrink as you learn more. But if the white balls are heavier and sink, your sample leans toward red every time. That is statistical bias, and more data may not help. More data ``can even consistently worsen a biased prediction.''

A cognitive bias works ``in an analogous way'': it is a systematic error in how we think. If the red balls could cure your brother's disease, you might overestimate how many there are because you wish they were many. Biases in yourself are harder to fix, but they are the place to start: ``if you can't trust your brain, how can you trust anything else?''

Most people guess that a shy stranger is a librarian rather than a salesperson, though salespeople are seventy-five times as common in the United States. \nb{This is a shortened version of Tversky and Kahneman's ``Steve'' problem (1974).} That is base rate neglect. People judge by how well the traits seem to fit together and ignore how common each kind of person is. Another bias is the sunk cost fallacy, the pull to stay committed to what you have already spent resources on when you should cut your losses.

Knowing about these biases does not make you immune to them. In one study, people predicted that knowing a painting was by a famous artist would make it harder for them to judge it neutrally. When tested, they showed exactly that bias, and afterward they claimed their judgments had been objective. People also see bias in others more easily than in themselves. When they look inward they find no ``biased-feeling thoughts,'' so they conclude that they are less biased than everyone else.

``Yet it is possible to recognize and overcome biases.'' My evidence is one sentence, without a source: people make fewer base-rate errors when they think in frequencies. The book's approach is to teach ``why good reasoning works.'' ``To the extent this volume does its job,'' I compare it to a finding by Sebastian Serfas. Accounting courses reduced people's sunk-cost bias, while years of work experience did not. Serfas calls the difference expertise, as opposed to experience. Expertise means understanding why the bias arises, spotting it in a given setting, and having tools to counter it. The book aims at that kind of expertise about human bias and self-deception. \nb{The book is a collection of blog essays, not a course, and the introduction gives no evidence that reading it has the courses' effect. Its method, understanding, also sits uneasily with the Preface just before it, which calls the largest mistake of the posts not seeing that the problem ``was figuring out how to practice it, not knowing the theory.''}

The book began as Yudkowsky's essays of 2006 to 2009. Its organizing idea is Alfred Korzybski's: ``The map is not the territory.'' A correct map is useful because its structure is similar to the territory's. The same holds for beliefs and words. The four sequences take this up in turn. ``Predictably Wrong'' is about the systematic ways our beliefs fail to map the world. ``Fake Beliefs'' asks what makes a belief a map at all. ``Noticing Confusion'' is about how the brain's mapping works, and ``Mysterious Answers'' about what happens when these points meet. A dialogue about truth itself, ``The Simple Truth,'' follows. Humans, in Dan Ariely's word, are ``predictably'' irrational, and understanding how we err can help us build a better future.
